\documentclass{article}
\usepackage{algorithm}
\usepackage{algpseudocode}
\usepackage{listings}
\usepackage{xcolor}
\usepackage{amsmath}


\lstset{
    language=Python,
    basicstyle=\ttfamily\small,
    keywordstyle=\color{blue},
    commentstyle=\color{green!60!black},
    stringstyle=\color{red},
    showstringspaces=false,
    frame=single,
    numbers=left,
    numberstyle=\tiny,
    breaklines=true
}
\usepackage[utf8]{inputenc}
\usepackage[a4paper,top=0.8in,bottom=1.0in,left=1.0in]{geometry}

\author{}
\date{}
\title{\textbf{Practical 4: Implementation and
complexity analysis of
algorithms based on dynamic
programming.
Floyd–Warshall Algorithm
(All-Pairs Shortest Paths)}}

\begin{document}
\maketitle

\vspace{0.1mm}
\begin{center}
    \LARGE\textbf{Floyd-Warshall Algorithm}\end{center}


{\normalsize
\section*{Objective}
\begin{itemize}
    \item Unlike single-source algorithms like Dijkstra or Bellman-Ford, it determines the optimal distance between all possible start and end nodes simultaneously.
    \item It can correctly compute shortest paths in graphs containing negative edge weights, provided the graph contains no negative cycles.
    \item t systematically improves path estimates by considering each vertex as a potential intermediate pivot to find shorter routes between other nodes. 
\end{itemize}

\section*{Theory}
\paragraph{} 
The Floyd-Warshall algorithm is a graph algorithm that is deployed to find the shortest path between all the vertices present in a weighted graph. This algorithm is different from other shortest path algorithms; to describe it simply, this algorithm uses each vertex in the graph as a pivot to check if it provides the shortest way to travel from one point to another.
\paragraph{}
There are mainly four steps in the algorithm:
\begin{itemize}
    \item Construct an adjacency matrix A with all the costs of edges present in the graph. If there is no path between two vertices, mark the value as $\infty$
    \item Derive another adjacency matrix $A_1$ from $A$ by keeping the first row and the first column of the original adjacency matrix unchanged. For the remaining entries, denoted by $A_1[i,j]$, update them according to the following rule:

\[
A_1[i,j] =
\begin{cases}
A[i,1] + A[1,j], & \text{if } A[i,j] > A[i,1] + A[1,j], \\
A[i,j], & \text{otherwise.}
\end{cases}
\]

In this step, the pivot vertex is the first vertex, i.e., $k = 1$.
    \item Repeat Step 2 for all the vertices in the graph by changing the k value for every pivot vertex until the final matrix is achieved.
    \item  The final adjacency matrix obtained is the final solution with all the shortest paths.
\end{itemize}
\pagebreak
\section*{Algorithm}
\begin{algorithm}
\caption{Quick Sort}
\begin{algorithmic}[1]

\Procedure{FloydWarshall}{$W, n$}
    \State $D \gets W$
    \For{$k \gets 0$ to $n-1$}
        \For{$i \gets 0$ to $n-1$}
            \For{$j \gets 0$ to $n-1$}
                \If{$D[i][k] + D[k][j] < D[i][j]$}
                    \State $D[i][j] \gets D[i][k] + D[k][j]$
                \EndIf
            \EndFor
        \EndFor
    \EndFor
    \State \Return $D$
\EndProcedure
\end{algorithmic}
\end{algorithm}

\section*{Python Code}
\begin{lstlisting}[={}]
def Floyd_warsall(w, n):
    D = [row[:] for row in w]

    for k in range(n):
        for i in range(n):
            for j in range(n):
                if D[i][k] + D[k][j] < D[i][j]:
                    D[i][j] = D[i][k] + D[k][j]
    return D
\end{lstlisting}

\subsection*{Output}
\begin{lstlisting}
INF = float('inf')
n = 4
W = [
    [0, INF, 3, INF],
    [2, 0, INF, INF],
    [INF, 7, 0, 1],
    [6, INF, INF, 0]
]

d = Floyd_warsall(W, n)
for i in range(n):
    for j in range(n):
        print(d[i][j], end='\t')
    print()

0       10      3       4
2       0       5       6
7       7       0       1
6       16      9       0
\end{lstlisting}
\pagebreak
\section*{Complexity Analysis}
Time Complexity:
\begin{itemize}
    \item $O(n^3)$
\end{itemize}
Space Complexity:
\begin{itemize}
    \item $O(n^2)$
\end{itemize}
}

\end{document}